\section{A canonical quadrilateral kernel for exact disc counts}
\label{sec:quadrilateral-core}

The weighted census answers a stronger question than an aggregate topology
summary, but it still receives endpoints whose bit lengths reflect every
normal-coordinate multiplicity. The maintained common-gcd reduction does not
remove a large additive collection of vertex-link surfaces. Such additions
can make the gcd of all $7t$ coordinates equal to one while almost every
coordinate has thousands of bits. For the particular question of counting
essential disc components, we can remove both sources of numerical growth.

The coordinate canonicalization is classical Q-theory. Tollefson introduced
quadrilateral coordinates, and Burton describes equal quadrilateral data,
vertex-link ambiguity, canonical representatives, positive homogeneity and
minimum-triangle subtraction explicitly
\cite{Tollefson1998,Burton2009}. We use these facts with proofs because their
exact hypotheses and their connection to the maintained interval interface
matter here. The delivered specialization is a source-bound essential-disc
count, together with its parameter-sensitive query bound and replayable
arithmetic reduction.

\subsection{Remove precisely the normal vertex-link components}

Let $x$ be an admissible normal vector in the supported finite triangulation
with $t$ tetrahedra. For a global vertex $v$, let $\mathcal C_v$ be the set of
tetrahedron corners identified with $v$, and write $n_v=|\mathcal C_v|$.
The sets $\mathcal C_v$ partition the $4t$ corners. Write $T_c(x)$ for the
triangle coordinate at corner $c$, and let $L_v$ be the vector with one
triangle at every corner in $\mathcal C_v$ and zero elsewhere. It represents
the vertex link: a sphere at an interior vertex and a disc at a boundary
vertex. Define
\[
 \ell_v=\min_{c\in\mathcal C_v}T_c(x),
 \qquad Y=x-\sum_v\ell_vL_v.
\]
Thus $Y$ has a zero triangle coordinate at some corner of every global
vertex. Its quadrilateral coordinates equal those of $x$.

\begin{lemma}[Geometric meaning of vertex-link subtraction]
\label{lem:vertex-link-peeling}
The vector $Y$ is admissible, and the embedded normal surface represented by
$x$ is normally isotopic to the disjoint union of the surface represented by
$Y$ and $\ell_v$ copies of the vertex link at every $v$.
\end{lemma}
\begin{proof}
Nonnegativity follows from the definition of each minimum, and no
quadrilateral coordinate changes. A matching equation across a face contains
one triangle coordinate from each of two corners identified with the same
global vertex. Subtracting $\ell_v$ from both preserves that equation.

Realize the admissible vector $Y$ as a normal surface. Since the surface is
finite and misses the triangulation vertices, sufficiently small vertex
neighborhoods are disjoint from it. Nested frontiers of these neighborhoods
give the requested vertex-link copies, also disjoint from one another.
Their disjoint union with the realization of $Y$ has coordinate vector $x$.
Uniqueness of a normal realization with its admissible coordinates, up to
normal isotopy, gives the asserted decomposition. This argument uses an
actual disjoint union; it does not assume an arbitrary normal sum preserves
its summands as components.
\end{proof}

Every boundary vertex-link disc has contractible boundary on the boundary
torus. Its presence therefore increases the number of discs without
increasing the number of \emph{compressing} discs. Interior vertex links
are closed and contribute no such disc either. These facts justify removing
the links before the specialized query, even when their multiplicities
dominate the input size.

\subsection{Quadrilateral content is the full content of the peeled vector}

Suppose at least one quadrilateral coordinate is nonzero, and let
\[
 g=\gcd\{Q_{\tau,j}(x):1\le\tau\le t,\ 1\le j\le3\}.
\]
Zero entries participate in the ordinary gcd convention. Let
$Q_{\max}$ be the largest quadrilateral coordinate and set $r=Q_{\max}/g$.
In particular $g\ge1$ and $r\ge1$.

\begin{theorem}[Canonical primitive quadrilateral core]
\label{thm:primitive-q-core}
Every coordinate of $Y$ is divisible by $g$, and its full coordinate gcd is
exactly $g$. The vector $P=Y/g$ is admissible, primitive and has a zero
triangle at every global vertex. Moreover,
\[
 T_c(P)\le(n_v-1)r\quad(c\in\mathcal C_v),
 \qquad \max_i P_i\le(4t-1)r.
\]
The vector $P$ is uniquely determined by the primitive quadrilateral vector
$Q(x)/g$. If all quadrilateral coordinates of $x$ are zero, then $Y=0$.
\end{theorem}
\begin{proof}
Consider the graph whose vertices are the corners in $\mathcal C_v$, with
edges recording their identifications across paired faces. It is connected
because those face identifications define the global vertex class. Along
each graph edge, the normal matching equation has the form
\[
 T_c(Y)-T_{c'}(Y)=Q'-Q,
\]
where $Q,Q'$ are individual quadrilateral coordinates, possibly zero.
The right side is divisible by $g$. Choose a corner with triangle
coordinate zero and follow a graph path to any other corner. Its triangle
coordinate is divisible by $g$ as well. All quadrilateral coordinates are
already divisible by $g$. Conversely, a common divisor of all coordinates
must divide their quadrilateral subset, whose gcd is $g$. The full gcd is
therefore exactly $g$.

After division, every quadrilateral appearing in one matching equation lies
between zero and $r$, so each triangle difference has absolute value at most
$r$. A simple path from a zero corner has at most $n_v-1$ edges. Summing the
differences proves the triangle bound. The global bound follows from
$n_v\le4t$ and the bound $r$ on quadrilateral coordinates.

If two vectors have the same quadrilateral coordinates, the difference
between their triangles is constant on each connected corner graph. Requiring
minimum zero at each vertex fixes every such constant. This proves uniqueness
of the vertex-reduced lift and hence of $P$. Finally, if every quadrilateral
vanishes, all triangle differences on each corner graph vanish; its minimum
zero forces every residual triangle to vanish.
\end{proof}

The theorem is invoked only after validating an actual supplied standard
normal vector. It does not assert that arbitrary quadrilateral data has a
finite normal lift, and it does not replace the matching and quadrilateral
constraints with an unchecked Q-coordinate input.

\subsection{Essential-disc counts survive the nonlinear sheet behavior}

Let $\D(z)$ denote the number of compressing-disc components in the embedded
normal surface of vector $z$. The scalar behavior of this number is simpler
than the scalar behavior of the full component count.

\begin{lemma}[Homogeneity of the essential-disc count]
\label{lem:disc-homogeneity}
For every admissible normal vector $z$ in the supported orientable manifold
and every positive integer $k$,
\[
 \D(kz)=k\D(z).
\]
\end{lemma}
\begin{proof}
Realize the surface with coordinates $kz$ in the normal interval
neighborhood of the surface with coordinates $z$. Local normal sheet orders
are either preserved or reversed. A two-sided component therefore produces
$k$ disjoint copies. Each compressing disc is two-sided, and these copies
have isotopic essential boundaries on the ambient torus.

A one-sided component has nontrivial normal monodromy: sheet $j$ is paired
with sheet $k-1-j$. For even $k$, this produces $k/2$ copies of its connected
normal orientation double. For odd $k$, it produces $(k-1)/2$ such doubles
and one original one-sided component. The ambient manifold is orientable,
so a one-sided surface is nonorientable. If it has boundary, its Euler
characteristic is at most zero, and the orientation double also has Euler
characteristic at most zero. Neither can be a disc. If the component is
closed, every such cover remains closed. Thus one-sided components never
create discs through this scaling process. Exactly the original compressing
discs contribute, each with multiplicity $k$.
\end{proof}

\begin{corollary}[Exact count on the primitive core]
\label{cor:core-disc-count}
With the notation above, if $g>0$ then
\[
 \boxed{\D(x)=g\D(P).}
\]
If all quadrilateral coordinates vanish, then $\D(x)=0$ and no orbit query
is needed.
\end{corollary}
\begin{proof}
Apply Lemma~\ref{lem:vertex-link-peeling}; none of the removed vertex links
contributes an essential disc. Then apply Lemma~\ref{lem:disc-homogeneity}
to $Y=gP$. In the triangle-only case the peeled vector is zero.
\end{proof}

The count-only interface deliberately reports no full component count for
$x$. For example, twice a normal M\"obius band is an annulus, so multiplying
its original component count by two is incorrect. Both surfaces have zero
compressing-disc components, consistently with the corollary. The separate
unreduced census remains the appropriate interface for complete component
histograms and coordinate extraction.

The corollary asserts scalar homogeneity after a specific geometric
decomposition. It does not assert additivity of $\D$ under arbitrary
compatible normal-vector addition: regular exchanges in a general normal
sum can change the connected components. Vertex links are removable here
because Lemma~\ref{lem:vertex-link-peeling} realizes that particular addition
as a disjoint union.

\subsection{Input arithmetic and reduced query complexity}

There are at most $t$ nonzero quadrilateral coordinates by admissibility.
The coefficient bound gives a convenient coarse disc-count estimate
\[
 W(P)=\sum_iP_i\le(16t^2-3t)r\le16t^2r.
\]
The maintained edge-point orbit universe has at most $4W(P)$ points, so its
endpoints use
\[
 b_0=O\bigl(\log(t+1)+\log(r+1)\bigr)
\]
bits. This bound is independent of $g$ and the vertex-link multiplicities.
Only one three-weight component query is required on $P$.

For an explicit accounting, let $B$ bound the original coordinate bit
lengths, let $V(t,B)$ denote full source validation and decoding, and let
$\mathcal P(t,b)$ denote the polynomial cost of the delivered three-weight
census on a source whose coordinates and endpoints have $O(b)$ bits.
Writing $G(B)$ and $\operatorname{Div}(B)$ for binary gcd and division costs,
the count-only route has the bound
\[
\begin{split}
 V(t,B)&+O\!\left(tB+tG(B)+t\operatorname{Div}(B)\right)\\
 &+\mathcal P\!\left(t,O(\log(t+1)+\log(r+1))\right)
 +O\!\left(M(B+\log(t+1))\right),
\end{split}
\]
where $M$ accounts for the final binary multiplication. Comparisons and
subtractions for the vertex minima are charged in $tB$. A second validation
of the small core, as performed by the current wrapper, is included in
$\mathcal P$. In the triangle-only case the census term is absent.

The certificate binds the full original source, records the vertex-link
multiplicities, the divisor, the explicit core and its independently checked
three-weight proof, and verifies the final product. Thus the outer arithmetic
fields still depend on the original input size. Large multiplicity is not
removed from the cost of reading, validating or binding that input; it is
removed from every endpoint and weight in the expensive core orbit trace.

\begin{example}[An arbitrarily large source with full coordinate gcd one]
Take a primitive, vertex-reduced meridian vector $D$ in a one-vertex layered
solid torus, with quadrilateral gcd one, and let $L$ be the boundary vertex
link. For any positive binary integer $g$, set
\[
 x_g=gD+(g+1)L.
\]
The quadrilateral coordinates have gcd $g$, while a triangle at a zero corner
of $D$ equals $g+1$. Hence the gcd of all coordinates is one. Common
multiplicity reduction alone leaves the original large endpoints in place.
The canonical kernel removes $g+1$ vertex-link discs, divides by $g$, and
recovers the fixed core $D$. It returns exactly $g$ compressing-disc
components. For a fixed triangulation and deterministic scheduler, its core
orbit trace is independent of the number of bits of $g$.
\end{example}

This is a concrete improvement for a supplied-vector subroutine. Primitive
quadrilateral coordinates may themselves remain large, and the operation
does not search for a useful normal vector or construct its knot-exterior
provenance. Those separate geometric obligations continue to control the
general recognition problem.
